% Part: incompleteness
% Chapter: arithmetization-syntax
% Section: coding-terms

\documentclass[../../../include/open-logic-section]{subfiles}

\begin{document}

\olfileid[tr]{inc}{art}{trm}
\olsection{Terimlerin Kodlanması}

\begin{explain}
Bir terim yalnızca belirli türde bir simge dizisidir: sabitlerden ve
değişkenlerden, terimlerin kuruluş kurallarına göre tümevarımlı biçimde
kurulur. Simgelere yönelik bir kodlama şeması ile sayı dizilerini kodlamanın
bir yolu kullanılarak simge dizileri sayılarla kodlanabildiğinden, terimlere
Gödel numaraları atamak zor değildir. Asıl güçlük, bir sayının düzgün
kurulmuş bir terimin Gödel numarası olması özelliğinin hesaplanabilir, hatta
ilkel özyinelemeli olduğunu göstermektir.
\end{explain}

!!^{variable}s ve !!{constant}s en basit terimlerdir; $x$'in böyle bir
terimin Gödel numarası olup olmadığını sınamak kolaydır: $x$, bir $i$ için
$\Gn{\Obj v_i}$ ise $\fn{Var}(x)$ geçerlidir. Başka bir deyişle $x$,
uzunluğu $1$ olan bir dizidir ve tek elemanı $(x)_0$, bir $\Obj v_i$
!!{variable}{}inin kodudur; yani $x$, bir $i$ için
$\tuple{\tuple{1, i}}$'dir. Benzer biçimde $x$, bir $i$ için
$\Gn{\Obj c_i}$ ise $\fn{Const}(x)$ geçerlidir. Böyle bir $i$ varsa
$i<x$ olmak zorunda olduğundan, bu bağıntıların ikisi de ilkel
özyinelemelidir:
\begin{align*}
  \fn{Var}(x) & \defiff \bexists{i<x}{x = \tuple{\tuple{1, i}}}\\
  \fn{Const}(x) & \defiff \bexists{i<x}{x = \tuple{\tuple{2, i}}}
\end{align*}

\begin{prop}
\ollabel{prop:term-primrec}
$x$'in sırasıyla bir terimin ya da kapalı bir terimin Gödel numarası olması
durumunda geçerli olan $\fn{Term}(x)$ ve $\fn{ClTerm}(x)$ bağıntıları ilkel
özyinelemelidir.
\end{prop}

\begin{proof}
Bir $s$ simge dizisi ancak ve ancak $s$ teriminin terimlerin kuruluş
kurallarına göre !!{constant}s{}den ve !!{variable}s{}den nasıl kurulduğunu
kaydeden bir $s_0$, \dots, $s_{k-1}=s$ terim dizisi varsa bir terimdir.
Böyle bir kuruluş dizisi kuruluş kurallarına uyuyorsa her $i<k$ için
aşağıdakilerden biri geçerlidir:
\begin{enumerate}
\item $s_i$, bir $\Obj v_j$ !!a{variable}{}dir; ya da
\item $s_i$, bir $\Obj c_j$ !!a{constant}{}tir; ya da
\item $s_i$, $i$ konumundan önce geçen $n$ tane $t_1$, \dots, $t_n$
  teriminden, $n$-li bir $\Obj f^n_j$ !!{function}{}u kullanılarak kurulmuştur.
\end{enumerate}
Gödel numaraları üzerindeki karşılık gelen bağıntının ilkel özyinelemeli
olduğunu göstermek için bu koşulu ilkel özyinelemeli, yani ilkel özyinelemeli
fonksiyonlar, bağıntılar ve sınırlı niceleme kullanarak ifade etmeliyiz.

$y$'nin $s_0$, \dots, $s_{k-1}$ dizisini kodlayan sayı olduğunu, yani
$y=\tuple{\Gn{s_0},\dots,\Gn{s_{k-1}}}$ olduğunu varsayalım. $y$, ancak ve
ancak her $i<k$ için aşağıdakilerden biri geçerliyse Gödel numarası $x$ olan
terimin bir kuruluş dizisini kodlar:
\begin{enumerate}
\item $\fn{Var}((y)_i)$; ya da
\item $\fn{Const}((y)_i)$; ya da
\item her $z_l$, bir $i'<i$ için $(y)_{i'}$'ye eşit olacak biçimde bir $n$
  ve $z=\tuple{z_1,\dots,z_n}$ sayısı vardır ve
\[
(y)_i = \Gn{\Obj f^n_j(} \concat \fn{flatten}(z) \concat \Gn{)},
\]
\end{enumerate}
ayrıca $(y)_{k-1}=x$ olur. ($\fn{flatten}(z)$ fonksiyonu
$\tuple{\Gn{t_1},\dots,\Gn{t_n}}$ dizisini $\Gn{t_1,\dots,t_n}$'ye
dönüştürür ve ilkel özyinelemelidir.)

(3)'teki $j$, $n$ indisleri, $t_l$ terimlerinin $z_l$ Gödel numaraları ve
$\tuple{z_1,\dots,z_n}$ dizisinin $z$ kodu $y$'den küçüktür. Yukarıdaki
$k$'yi $\len{y}$ ile değiştirebiliriz. Dolayısıyla ``$y$, Gödel numarası $x$
olan terimin bir kuruluş dizisinin kodudur'' ifadesini bu bağıntının ilkel
özyinelemeli olduğunu gösterecek biçimde yazabiliriz.

Şimdi yalnızca $y$ üzerinde ilkel özyinelemeli bir sınır bulunduğuna ikna
olmalıyız. $x$ bir terimin Gödel numarasıysa, en çok $\len{x}$ terimli bir
kuruluş dizisi vardır; çünkü $s$'nin kuruluş dizisindeki her terim $s$
içinde bir konumda başlamalıdır ve iki alt terim aynı konumda başlayamaz.
$s$'nin her alt teriminin Gödel numarası elbette $\le x$'tir. Dolayısıyla
$k=\len{x}$ olmak üzere kodu $\le p_{k-1}^{k(x+1)}$ olan bir kuruluş dizisi
her zaman vardır.

$\fn{ClTerm}$ için !!{variable}s{}i kapsayan koşulu çıkarmak yeterlidir.
\end{proof}

\begin{prob}
$\tuple{\Gn{t_1},\dots,\Gn{t_n}}$ dizisini $\Gn{t_1,\dots,t_n}$'ye
dönüştüren $\fn{flatten}(z)$ fonksiyonunun ilkel özyinelemeli olduğunu
gösterin.
\end{prob}

\begin{prop}\ollabel{prop:num-primrec}
  $\fn{num}(n)=\Gn{\num{n}}$ fonksiyonu ilkel özyinelemelidir.
\end{prop}

\begin{proof}
  $\fn{num}(n)$'yi ilkel özyinelemeyle tanımlarız:
  \begin{align*}
    \fn{num}(0) & = \Gn{\Obj 0}\\
    \fn{num}(n+1) & = \Gn{\prime(} \concat \fn{num}(n) \concat \Gn{)}.
  \end{align*}
\end{proof}

\end{document}
